\documentclass[11pt,a4paper]{article}
\usepackage[T1]{fontenc}
\usepackage{lmodern,microtype}
\usepackage[margin=27mm]{geometry}
\usepackage{amsmath,amssymb,amsthm,mathtools}
\usepackage{booktabs,longtable,array,enumitem,needspace}
\usepackage[dvipsnames]{xcolor}
\usepackage{listings}
\usepackage{hyperref}
\hypersetup{colorlinks=true,linkcolor=MidnightBlue,urlcolor=MidnightBlue,
 citecolor=MidnightBlue,pdftitle={Robin's theorem in iPRA: finite components and analytic obligations},
 pdfauthor={Research prepared for Vladimir Reshetnikov}}
\setlength{\emergencystretch}{2em}
\setlength{\parskip}{4pt}
\setlength{\parindent}{0pt}
\newtheorem{theorem}{Theorem}[section]
\newtheorem{lemma}[theorem]{Lemma}
\newtheorem{proposition}[theorem]{Proposition}
\newtheorem{corollary}[theorem]{Corollary}
\theoremstyle{definition}
\newtheorem{definition}[theorem]{Definition}
\newtheorem{remark}[theorem]{Remark}
\newcommand{\ipra}{\mathsf{iPRA}}
\newcommand{\pra}{\mathsf{PRA}}
\newcommand{\RH}{\mathsf{RH}}
\newcommand{\RR}{\mathbb R}
\newcommand{\CC}{\mathbb C}
\newcommand{\NN}{\mathbb N}
\newcommand{\QQ}{\mathbb Q}
\newcommand{\Pos}{\mathsf{Pos}}
\newcommand{\Neg}{\mathsf{Neg}}
\newcommand{\Bad}{\mathsf{Bad}}
\newcommand{\Rstr}{\mathsf{Robin}_{<}}
\newcommand{\Rweak}{\mathsf{Robin}_{\geq}}
\newcommand{\Rdn}{\mathsf{Robin}_{\neg\neg}}
\newcommand{\RHc}{\mathsf{RH}_{\rm cert}}
\newcommand{\code}[1]{\texttt{\small #1}}
\lstset{basicstyle=\ttfamily\footnotesize,breaklines=true,
 columns=fullflexible,keepspaces=true,frame=single,
 rulecolor=\color{black!20},backgroundcolor=\color{black!2}}
\begin{document}
\begin{titlepage}
\centering
\vspace*{12mm}
{\Large Constructive number theory and proof foundations\par}
\vspace{14mm}
{\huge\bfseries Robin's theorem in intuitionistic\\[3mm]
primitive recursive arithmetic\par}
\vspace{8mm}
{\Large Verified finite components and a precise analytic completion problem\par}
\vspace{12mm}
{\large Prepared for Vladimir Reshetnikov\\19 November 2026\par}
\vfill
\begin{minipage}{0.92\textwidth}
\textbf{Status of the requested theorem.}
This report does \emph{not} establish the full equivalence in $\ipra$.
It gives complete human proofs of explicit finite arithmetic reductions,
an exact interval construction, a finite verification through $100000$,
and a locally compiled constructive Lean core. It identifies the analytic
and object-theory obligations which remain before the requested equivalence
can honestly be claimed.

The distinction is substantive: neither a classical Lean theorem nor a
record whose fields request the missing proofs constitutes a proof in
first-order $\ipra$. No impossibility or independence result is asserted.
\end{minipage}
\vfill
{\small Source snapshots, exact certificate data, replay code, Lean source,
and verification records accompany this report.}
\end{titlepage}
\tableofcontents
\newpage

\section{The exact problem and the results obtained}

Robin's theorem is the classical equivalence
\begin{equation}\label{eq:target}
 \RH\quad\Longleftrightarrow\quad
 \forall n\ge5041,\quad \sigma(n)<e^\gamma n\log\log n,
\end{equation}
where
\[
 \sigma(n)=\sum_{\substack{1\le d\le n\\d\mid n}}d,\qquad
 H_m=\sum_{j=1}^{m}\frac1j,\qquad
 \gamma=\lim_{m\to\infty}(H_m-\log m).
\]
The requested task is to prove the equivalence, not to prove $\RH$.
The numerical cutoff and the strict inequality are both essential.
The statement agrees with the supplied MathWorld page and Robin's original
Theorem~1 \cite{MathWorld,Robin}.

For clarity put
\begin{equation}\label{eq:gap}
 B(n)=e^\gamma\log\log n,\qquad
 A(n)=\frac{\sigma(n)}n,\qquad
 G(n)=B(n)-A(n).
\end{equation}
The right side of \eqref{eq:target} is $G(n)>0$ for every $n>5040$.
A positive real inequality is represented below by an explicit positive
rational lower bound, not merely by failure to certify a negative value.

\subsection{What is proved here}

The main completed mathematical components are as follows.
\begin{enumerate}[leftmargin=*]
\item Explicit rational enclosures for logarithms, $\gamma$, $e^\gamma$,
and $B(n)$, with primitive-recursive precision schedules. The correctness
proof uses finite geometric identities, finite telescoping, and specified
tails.
\item Exact certificates for all $99999$ integers $2\le n\le100000$.
Every comparison is strict; no case is left undetermined.
\item A uniform infinite subcase: every $n\ge5041$ with at most three
distinct prime factors satisfies Robin's inequality.
\item An elementary counterexample transfer. For $m\ge16$ and
\[
 F(m^2)=e^\gamma\log\vartheta(m^2)
        \prod_{p\le m^2}(1-p^{-1}),
\]
the inequality $-\log F(m^2)>12/m$ produces a specific counterexample
$\operatorname{lcm}(1,\ldots,m^2)>5040$.
\item A constructive Lean core for the logic of strict certificates,
negative certificates, stable negative RH, and bounded witness search.
Eight printed axiom audits return no axioms.
\item A useful conditional metatheorem: an appropriately arithmetized
\emph{classical PRA} proof of this $\Pi^0_1$--$\Pi^0_2$ equivalence would
yield an intuitionistic PRA proof. Thus the use of classical logic by
itself is not an impossibility argument.
\end{enumerate}

\begin{center}
\small
\begin{tabular}{@{}>{\raggedright\arraybackslash}p{.34\textwidth}
 >{\raggedright\arraybackslash}p{.58\textwidth}@{}}
\toprule
Component & Status in this archive\\
\midrule
Finite arithmetic and LCM arguments & Complete human proofs; object iPRA
derivations outlined, not constructed.\\
Finite interval data & Generated and replayed using exact integers.\\
Constructive logical shell & Compiled locally in Lean 4.32.0; no axioms
in the eight audited declarations.\\
Concrete zero-certificate checker & Detailed arithmetic specification;
not implemented or proved as an iPRA derivation.\\
RH-dependent sharp analytic bound & Identified precisely from a classical
development; iPRA interpretation not established here.\\
Nicolas--Landau analytic converse & Its finite output reduction is proved;
the required iPRA analytic argument remains open in this work.\\
Full equivalence \eqref{eq:target} in iPRA & Not established.\\
\bottomrule
\end{tabular}
\end{center}

\subsection{Why the distinction matters}

A computable real function can be evaluated at every specified precision
without admitting a primitive-recursive bound on a search for a strictly
positive margin. Likewise, a classical theorem in Lean's ambient logic
does not identify the induction and comprehension used by its proof.
Conversely, the existence of classical steps does not establish that the
theorem is unprovable constructively. The remaining mathematical task is
to supply the appropriate arithmetic proof, or a proof in a conservative
extension with an applicable conservation theorem.

This report therefore provides useful proved components and a precise
completion route, while keeping the uncompleted analytic theorem visible.

\section{The object theory and its arithmetic representations}

\subsection{Language, axioms, and inference}

We use first-order intuitionistic primitive recursive arithmetic. Its
objects are natural numbers. There is a function symbol for each fixed
primitive-recursive definition, formed from zero, successor, projections,
composition, and primitive recursion. Its nonlogical axioms are the
defining equations, successor nonzero and injectivity, and the induction
scheme
\[
 \bigl(Q(0,\vec a)\wedge
 \forall t(Q(t,\vec a)\to Q(t+1,\vec a))\bigr)
 \to\forall t\,Q(t,\vec a)
\]
for quantifier-free $Q$. Logic is intuitionistic first-order natural
deduction with equality; negation means implication to false. No excluded
middle, Markov principle, countable choice, or unrestricted induction is
added.

Each fixed bounded predicate has a primitive-recursive Boolean
characteristic function. A bounded universal assertion is compiled into
the Boolean fold which checks each entry. Induction on the value of this
fold is quantifier-free. This observation must be applied to a
\emph{specified} predicate and specified bounds; it is not a general
reflection principle for arbitrary statements.

\subsection{Finite arithmetic codes}

Signed integers can be represented by a sign and a magnitude; rational
numbers by a signed numerator and a positive denominator. Reduction to
lowest terms is optional. Equality and order are checked by signed
cross-multiplication. Finite sequences use any standard primitive-recursive
pairing and list code with bounded projections. All constructions below
can instead use denominator-cleared integers.

For $n>0$, divisibility is checked by bounded division. Trial division
through $n$ gives a factorization trace. The prime predicate
\[
 \operatorname{Prime}(p)\;\equiv\;
 p\ge2\ \wedge\
 \forall d\le p\,\bigl(d\mid p\to(d=1\vee d=p)\bigr)
\]
is bounded and hence primitive recursive. The divisor sum above is a
bounded sum. Factorials, binomial coefficients, finite products, integer
square roots, and
\[
 L_1=1,\qquad L_{j+1}=\operatorname{lcm}(L_j,j+1)
\]
are primitive recursive. Euclid's gcd algorithm may be replaced by
bounded search if a particularly simple totality argument is wanted.

Exact rational intermediate values can be large. Their size does not
invalidate primitive recursiveness: every loop and every recursion in
these definitions has an explicitly supplied natural-number argument.
Claims about computational practicality are kept separate.

\subsection{Real names are abbreviations for fixed programs}

A rational interval name $x$ consists, externally, of a fixed pair of
primitive-recursive endpoint functions $l_x(k),u_x(k)$ and proofs of
compatibility and an explicit shrinking bound. Its rational comparisons
are arithmetic predicates. No quantification over arbitrary real
numbers is needed for the numerical Robin side.

We write $x>0$ positively as
\[
 \exists k\quad l_x(k)>0.
\]
The negative weak statement $x\ge0$ is
$\forall k\,\neg(u_x(k)<0)$. These are different formulas.
When different names denote the same quantity, a finite comparison
estimate must justify substitution; a semantic assertion in $\RR$ is
not automatically an object derivation.

There is no total primitive-recursive universal evaluator for all
primitive-recursive programs. The present convention fixes each program
before its representing term is introduced. The Lean shell later
parametrizes Boolean functions for readability, and explicitly does not
claim that arbitrary Lean functions are primitive recursive.

\section{Explicit interval arithmetic and its human correctness proof}
\label{sec:intervals}

\subsection{A logarithm enclosure from a finite geometric identity}

For $0\le z<1$ set
\[
 S_K(z)=2\sum_{j=0}^{K-1}\frac{z^{2j+1}}{2j+1}.
\]
The finite identity
\[
 \frac1{1-t^2}=\sum_{j=0}^{K-1}t^{2j}
                   +\frac{t^{2K}}{1-t^2}
\]
gives, on integrating between $0$ and $z$,
\begin{equation}\label{eq:logtail}
 0\le
 \log\frac{1+z}{1-z}-S_K(z)
 \le\frac{2z^{2K+1}}{(2K+1)(1-z^2)}.
\end{equation}
The upper estimate follows by replacing $(1-t^2)^{-1}$ by
$(1-z^2)^{-1}$ in the remainder. In a rational-name presentation the same
formula is the consistency bound between finite polynomials: subtract two
partial sums and bound the resulting finite tail by a geometric sum.
There is no unbounded search in either presentation.

For $1\le q\le2$, put $z=(q-1)/(q+1)\le1/3$. Thus
\begin{equation}\label{eq:uniform-log-tail}
 0\le\log q-S_K(z)
 \le r_K:=\frac{9}{4(2K+1)3^{2K+1}}
 \le3^{-2K}.
\end{equation}
For an arbitrary positive rational $q\ge1$, bounded binary range
reduction produces $q=2^a y$, with $a\ge0$ and $1\le y<2$.
The enclosure is the sum of $a$ copies of the enclosure of $\log2$ and
the enclosure of $\log y$. For $q<1$ evaluate $1/q$ and negate the
interval. Taking
\[
 K=k+a+4
\]
makes the interval width at most
$(a+1)3^{-2K}\le2^{-k}$. The inequality follows, for example, from
$a+1\le2^{a+1}$ and $9^K\ge2^{2K}$.

This defines a fixed primitive-recursive rational interval
$\operatorname{LogI}(q,k)$ of width at most $2^{-k}$. Powers,
denominators, and the finite summation length are explicit.

\subsection{Euler's constant without an unspecified convergence rate}

Let
\[
 a_m=H_m-\log m-\frac1{2m}\qquad(m\ge1).
\]

\begin{lemma}\label{lem:gamma}
For every $m\ge1$,
\begin{equation}\label{eq:gamma-bound}
 a_m\le\gamma\le a_m+\frac1{12m^2}.
\end{equation}
Moreover these bounds follow from a finite telescoping estimate with an
explicit tail.
\end{lemma}
\begin{proof}
Set $z=1/(2m+1)$. Direct subtraction gives
\[
 a_{m+1}-a_m
 =\frac1{2m}+\frac1{2(m+1)}-\log(1+1/m)
 =\frac{2z}{1-z^2}-2\operatorname{artanh}z.
\]
Using the geometric expansion already bounded in
\eqref{eq:logtail}, this difference is
\[
 4\sum_{j=1}^{\infty}\frac{j}{2j+1}z^{2j+1}.
\]
It is positive. Since $(2j+1)^{-1}\le1/3$ for $j\ge1$,
\begin{align*}
 0<a_{m+1}-a_m
 &\le\frac43\sum_{j=1}^{\infty}jz^{2j+1}
 =\frac{4z^3}{3(1-z^2)^2}\\
 &=\frac1{12m^2}-\frac1{12(m+1)^2}.
\end{align*}
Each displayed infinite sum can be replaced by a finite sum and its
explicit geometric tail. For $M>m$, finite telescoping therefore gives
\[
 0\le a_M-a_m
 \le\frac1{12m^2}-\frac1{12M^2}.
\]
This is a Cauchy estimate with an explicit modulus. It defines a limit
name; and $H_m-\log m=a_m+1/(2m)$ identifies that name with the usual
Euler constant. Letting the explicitly bounded remainder tend to zero
gives \eqref{eq:gamma-bound}. Equivalently one works directly with the
nested intervals, without a completed real field inside $\ipra$.
\end{proof}

The case $m=1$ gives $1/2\le\gamma\le7/12$, useful for every subsequent
domain check. For a rational enclosure at precision $k$, choose
$m=2^{k+4}$, evaluate $\log m$ with precision $k+4$, and use
\eqref{eq:gamma-bound}. The resulting width is less than $2^{-k}$.
Intersecting with $[1/2,7/12]$ is legitimate and ensures that all
exponential evaluations below lie in $[0,1]$.

\subsection{The exponential enclosure}

For $0\le x\le1$ put $E_K(x)=\sum_{j=0}^{K}x^j/j!$. The ratio of
consecutive omitted terms is at most $1/(K+2)$, so
\begin{equation}\label{eq:exptail}
 0\le e^x-E_K(x)\le\frac2{(K+1)!}.
\end{equation}
This too is a finite tail estimate, with the usual power series as its
external interpretation. If $x\in[l,u]\subseteq[0,1]$, monotonicity gives
\[
 E_K(l)\le e^x\le E_K(u)+\frac2{(K+1)!}.
\]
The difference $e^u-e^l$ is at most $3(u-l)$, since $e<3$.
The same Lipschitz bound follows by factoring $u^j-l^j$ in the finite
series and applying \eqref{eq:exptail}.

Use the $\gamma$ enclosure at precision $k+8$, and take $K=2k+32$.
Since $(K+1)!\ge2^K$, the width of the resulting enclosure of $e^\gamma$
is below $2^{-k}$. Intersecting with $[1,3]$ is harmless.

\subsection{Composing the double logarithm}

For $n\ge2$ let $\ell(n)$ denote the binary bit length of $n$ and put
\[
 r=k+\ell(n)+16.
\]
Enclose $\log n$ at precision $r$, then enclose the logarithms of its
two endpoints at the same precision. The first interval lies above
$1/2$, because $\log2\ge2/3$ and its width is tiny. On this domain
the logarithm has Lipschitz constant $2$. The double-log interval
therefore has width at most $4\cdot2^{-r}$. It can be intersected with
$[-1,n]$, which contains $\log\log n$ for $n\ge2$.

Multiply it by the $[1,3]$-bounded enclosure of $e^\gamma$ at precision
$r$, using the minimum and maximum of all four endpoint products.
The resulting rational interval $[L(n,k),U(n,k)]$ contains $B(n)$
and has width at most
\[
 (n+12)2^{-r}\le2^{-k}.
\]
This last inequality follows from $n<2^{\ell(n)}$ and the spare factor
$2^{16}$. All intermediate domains and all loop lengths have now been
specified. Define the fixed arithmetic predicates
\begin{align}
 \Pos(n,k)&:\Longleftrightarrow n\ge2\ \wedge\
                 nL(n,k)>\sigma(n),\label{eq:pos}\\
 \Neg(n,k)&:\Longleftrightarrow n\ge2\ \wedge\
                 nU(n,k)<\sigma(n).\label{eq:neg}
\end{align}
After clearing positive denominators these are primitive-recursive
Boolean checks on integers. They are mutually incompatible, even when
different precisions are used.

\subsection{The executable dyadic implementation}

The certificate generator uses the same inequalities with directed
rounding to denominator $S=2^{72}$, rather than the slower exact-rational
precision schedule just specified. It uses $32$ logarithm and exponential
terms and $m=2048$ for the Euler-constant bound. A rational operation
whose result is $a/b$ is rounded down by $\lfloor aS/b\rfloor$ for a
lower endpoint and up by $\lceil aS/b\rceil$ for an upper endpoint.
The four-corner product rule handles the negative value $\log\log2$.

For $\log q$ on $[1,2]$, the implementation rounds $z$, $z^2$, each odd
power, and each summand outward, then adds the ceiling of $Sr_K$.
For $\gamma$, it rounds $H_m$, $\log m$, $1/(2m)$ and $1/(12m^2)$
outward. For $e^\gamma$, it uses the nonnegative Taylor recurrence
$t_j=t_{j-1}x/j$ and adds the ceiling of $2S/(K+1)!$.
Monotonicity proves each loop invariant by induction on its bounded
counter. Final acceptance is solely one of
\[
 nL_{\rm int}>\sigma(n)S,\qquad
 nU_{\rm int}<\sigma(n)S.
\]

The dyadic program is a finite numerical verification tool. It is not
described as a formal proof checker for the entire object theory.
The human interval proof explains its exact certificate meanings;
the file \code{verification/arithmetic\_audit.md} records a separate
implementation audit.
Every accepted dyadic row supplies a positive rational gap. Choosing
$k$ with $2^{-k}$ smaller than that supplied gap converts it to a
canonical certificate \eqref{eq:pos} or \eqref{eq:neg}. A bound follows
immediately from the numerator and denominator bit lengths, so this
conversion does not require unbounded precision search.

\section{Finite verification and an infinite elementary subcase}

\subsection{Results through one hundred thousand}

The sieve computes $\sigma(n)$ by adding $d$ to every multiple of $d$.
Its invariant is that, after processing $d$, each array entry contains
the sum of its divisors at most $d$. Thus it is an exact bounded
arithmetic computation.

The generator classified every $2\le n\le100000$. The complete set of
strict failures in this range is
\begin{align*}
 &2,3,4,5,6,8,9,10,12,16,18,20,24,30,36,48,60,\\
 &72,84,120,180,240,360,720,840,2520,5040.
\end{align*}
Consequently every $5041\le n\le100000$ passes. This is a finite theorem
only. It does not justify any statement about all larger integers.

Here are the common exact enclosures, all with denominator
$S=4722366482869645213696$:
\begin{center}
\small
\begin{tabular}{@{}lll@{}}
\toprule
Quantity & Lower numerator & Upper numerator\\
\midrule
$\gamma$ & $2725823815493325450478$ & $2725823909318317689830$\\
$e^\gamma$ & $8410876523171405724439$ & $8410876690280509860046$\\
$B(5041)$ & $18024888747257598183496$ & $18024889105379965383064$\\
\bottomrule
\end{tabular}
\end{center}
For orientation, the last interval lies between $3.81691865$ and
$3.81691874$. These decimal figures are explanatory; the comparisons
in the archive use the integers.

Among the checked $n>5040$, the smallest certified lower bound for
$G(n)$ occurs at $n=10080$ and equals
\[
 \frac{2670582172179503638522368}
      {47601454147326023754055680}
 >0.05610.
\]
The CSV records $n,\sigma(n)$, both interval numerators, the binary
scale, and the accepted sign. Its SHA-256 is recorded in
\code{certificates/summary.json}. The independent replay recomputes
the arithmetic and checks row continuity, the advertised range, signs,
and the file digest.

\subsection{All integers with at most three distinct prime factors}

\begin{theorem}\label{thm:three}
If $n\ge5041$ has at most three distinct prime factors, then
$\sigma(n)<e^\gamma n\log\log n$.
\end{theorem}
\begin{proof}
For $n>1$ write $n=\prod_{i=1}^r p_i^{a_i}$ with distinct increasing
primes and $a_i\ge1$. Finite geometric summation gives
\[
 A(n)=\prod_{i=1}^r
       \frac{1-p_i^{-a_i-1}}{1-p_i^{-1}}
 <\prod_{i=1}^r\frac{p_i}{p_i-1}.
\]
Every factor removed from the numerator lies strictly between zero and
one. Since $p_1\ge2$, $p_2\ge3$, $p_3\ge5$ and $p/(p-1)$ decreases,
for $r\le3$ the right side is at most
\[
 2\cdot\frac32\cdot\frac54=\frac{15}4.
\]
The certificate above proves $B(5041)>15/4$ by a single integer
cross-multiplication. Both logarithms and $e^\gamma$ are increasing
or positive on the relevant domains, so $B(n)\ge B(5041)$.
Combining the two estimates proves the strict inequality.
\end{proof}

Only fixed finite factorization, finite product induction, elementary
logarithm monotonicity, and one displayed rational certificate are used.
This gives an actual infinite family of Robin inequalities without RH.
The object-theory proof would compile these specific finite recurrences;
it would not appeal to the unproved full criterion.

% Original finite bridge prepared for the Robin/iPRA report.
% This is a fragment intended for inclusion in the parent's document.
\section{An explicit finite least-common-multiple bridge}

This section proves an elementary reduction which does not use the Riemann
hypothesis or an oscillation theorem. Its input is one strict inequality at
a specified finite frontier. Its output is one explicitly constructed integer
violating Robin's inequality. All arithmetic sums and products are bounded.

For an integer $P\ge2$, define
\[
 \vartheta(P)=\sum_{\substack{p\le P\\p\text{ prime}}}\log p,
 \qquad
 a_p(P)=\max\{a\le P:p^a\le P\},
 \qquad
 \psi(P)=\sum_{\substack{p\le P\\p\text{ prime}}}a_p(P)\log p.
\]
The maximum defining $a_p(P)$ is a bounded search. Put
\[
 L_P=\operatorname{lcm}(1,\ldots,P),\qquad
 A(n)=\frac{\sigma(n)}n,
 \qquad
 M(n)=\log A(n)-\gamma-\log\log\log n.
\]
The last expression is used only for $n>e$. Since the logarithm is strictly
increasing, $M(n)>0$ is equivalent to
$\sigma(n)>e^\gamma n\log\log n$.

For $\vartheta(P)>1$, define
\[
 F(P)=e^\gamma\log\vartheta(P)
       \prod_{\substack{p\le P\\p\text{ prime}}}\left(1-\frac1p\right),
\]
and the two nonnegative losses
\[
 D(P)=\log\log\psi(P)-\log\log\vartheta(P),
 \qquad
 T(P)=\sum_{\substack{p\le P\\p\text{ prime}}}
       -\log\left(1-p^{-a_p(P)-1}\right).
\]

\begin{lemma}[Exact finite identity]
If $\vartheta(P)>1$, then
\[
 M(L_P)=-\log F(P)-D(P)-T(P).
\]
\end{lemma}
\begin{proof}
Unique prime factorization gives
$L_P=\prod_{p\le P}p^{a_p(P)}$, and hence $\log L_P=\psi(P)$.
The finite geometric-series formula and multiplicativity of $\sigma$ give
\[
 A(L_P)=\prod_{p\le P}
        \frac{1-p^{-a_p(P)-1}}{1-p^{-1}}.
\]
Taking logarithms yields
$\log A(L_P)=-T(P)-\sum_{p\le P}\log(1-p^{-1})$.
On the other hand,
\[
 \log F(P)=\gamma+\log\log\vartheta(P)
              +\sum_{p\le P}\log(1-p^{-1}).
\]
Substituting these equalities into the definition of $M(L_P)$ proves the
claim. Every logarithm is in its positive domain: $\psi(P)\ge\vartheta(P)>1$.
\end{proof}

\begin{lemma}[Losses at a square frontier]
Let $m\ge4$ be an integer, put $P=m^2$, and suppose
$\vartheta(P)\ge P/4$. Then
\[
 0\le D(P)\le\frac8m,\qquad
 0\le T(P)\le\frac4m,
 \qquad
 D(P)+T(P)\le\frac{12}m.
\]
\end{lemma}
\begin{proof}
Only primes $p\le m$ contribute to $\psi(P)-\vartheta(P)$.
For each such prime, $(a_p(P)-1)\log p\le\log P$; there are at most
$m$ such primes. Therefore
\[
 0\le\psi(P)-\vartheta(P)\le m\log P.
\]
For $v\ge u>1$, two applications of $\log(1+t)\le t$ give
\[
 \log\log v-\log\log u
 \le\frac{\log v-\log u}{\log u}
 \le\frac{v-u}{u\log u}.
\]
Since $P\ge16$,
$\log(P/4)\ge\tfrac12\log P$. Consequently
\[
 D(P)\le\frac{m\log P}{(P/4)\log(P/4)}\le\frac8m.
\]

For $0\le u\le1/2$, the geometric-series estimate gives
$-\log(1-u)\le u/(1-u)\le2u$.
If $p\le m$, maximality of $a_p(P)$ implies $p^{a_p(P)+1}>P$, so
$p^{-a_p(P)-1}<1/P$. These at most $m$ summands contribute at most $2/m$.
If $m<p\le P$, then $a_p(P)=1$, so the remaining summands contribute at most
\[
 2\sum_{j=m+1}^{P}\frac1{j^2}
 \le2\sum_{j=m+1}^{P}\frac1{j(j-1)}
 =2\left(\frac1m-\frac1P\right)\le\frac2m.
\]
Here the inequality is termwise and the equality is a finite telescoping
sum. This proves $T(P)\le4/m$, and addition proves the last claim.
\end{proof}

\begin{lemma}[A rough elementary Chebyshev lower bound]
For every integer $m\ge16$,
\[
 \vartheta(m^2)\ge\frac{m^2}4.
\]
\end{lemma}
\begin{proof}
Put $P=m^2$ and $h=\lfloor P/2\rfloor$.
Legendre's formula shows that the exponent of a prime $p$ in
$\binom{2h}{h}$ is
\[
 \sum_{j\ge1}\left(\left\lfloor\frac{2h}{p^j}\right\rfloor
       -2\left\lfloor\frac h{p^j}\right\rfloor\right).
\]
The sum is bounded by $j\le2h$. Every summand is either zero or one,
and it vanishes for $p^j>2h$. Thus the exponent is at most
$a_p(2h)$, proving $\binom{2h}{h}\mid L_{2h}$.
The central coefficient is the largest of the $2h+1$ coefficients in
$(1+1)^{2h}$; hence
\[
 \psi(P)\ge\log\binom{2h}{h}
 \ge2h\log2-\log(2h+1)
 \ge(P-1)\log2-\log(P+1).
\]

The logarithm series at $1/3$ gives $2/3\le\log2\le3/4$.
Also, since $m/16\ge1$ and $\log t\le t-1$,
\[
 \log m=4\log2+\log(m/16)
 \le3+(m/16-1)=2+m/16\le3m/16.
\]
Therefore $\log P=2\log m\le3m/8$, while
\[
 \log(P+1)\le\log(2P)\le\frac34+\frac{3m}8.
\]
The previous prime-power count gives
$\psi(P)-\vartheta(P)\le m\log P\le3m^2/8$. Combining the estimates,
\[
 \vartheta(m^2)
 \ge\frac{7m^2}{24}-\frac{3m}8-\frac{17}{12}
 =\frac{m^2}4+\frac{m^2-9m-34}{24}.
\]
Finally $m^2-9m-34=(m-16)(m+7)+78\ge0$ for $m\ge16$.
\end{proof}

\begin{theorem}[Explicit finite counterexample transfer]
Let $m\ge16$ be an integer. If
\[
 -\log F(m^2)>\frac{12}m,
\]
then the explicitly computed integer $N=L_{m^2}$ satisfies $N>5040$ and
\[
 \sigma(N)>e^\gamma N\log\log N.
\]
\end{theorem}
\begin{proof}
The previous two lemmas and the exact identity give
\[
 M(N)=-\log F(m^2)-D(m^2)-T(m^2)>0.
\]
Moreover $m^2\ge16$ and $L_{16}=720720$ divides $L_{m^2}$, so $N>5040$.
The equivalence between $M(N)>0$ and the displayed strict inequality was
already established.
\end{proof}

\begin{lemma}[Passing to a square frontier]
Let $Q\ge256$ be an integer, $m=\lfloor\sqrt Q\rfloor$, and $P=m^2$.
Then
\[
 \log F(P)\le\log F(Q)+\frac4m.
\]
Consequently $-\log F(Q)>16/m$ yields a Robin counterexample at $L_P$.
\end{lemma}
\begin{proof}
We have $m\ge16$ and $0\le Q-P\le2m$.
The term $\log\log\vartheta$ is nondecreasing, so
\[
 \log F(P)-\log F(Q)
 \le\sum_{P<p\le Q}-\log(1-p^{-1})
 \le\frac{2(Q-P)}{P+1}\le\frac4m.
\]
The previous theorem now applies.
\end{proof}

This lemma makes the connection with the conventional nonsquare oscillation
statement explicit. Suppose $c>0$, $0<b<1/2$, and an integer $Q\ge256$
satisfies $-\log F(Q)\ge cQ^{-b}$. Set $\delta=1/2-b$.
Since $m\ge\sqrt Q/2$, the additional condition $cQ^\delta>32$ implies
$-\log F(Q)>16/m$ and hence a counterexample. For rational $c,b$, choose
integers $B\ge\max(2,64/c)$ and $k\ge1/\delta$. Then the sufficient
cutoff $Q\ge\max(256,B^k)$ is primitive recursive in their rational codes.
Classically, Nicolas's excursions occur beyond every cutoff, so this
finishes the transfer from that analytic input. It does not itself prove
the analytic input or its iPRA interpretation.

The result also holds for $m\ge4$ when $\vartheta(m^2)\ge m^2/4$ is supplied
as an explicit hypothesis. The choice $m\ge16$ removes that hypothesis by
the completely elementary proof above. A supplied rational positive
gap determines a primitive-recursive precision for the log enclosure
and therefore a finite certificate. Strictness by itself is not being
used as an unbounded-search termination bound. The proof contains no
search for the input frontier.

\paragraph{What this does and does not finish.}
An analytic oscillation theorem would still have to supply an appropriate
finite frontier. Classical Nicolas--Landau oscillation gives such excursions
arbitrarily far out, of size $cP^{-b}$ for some $0<b<1/2$. The finite loss
is only of order $P^{-1/2}$, so that implication follows classically. To
obtain an iPRA theorem, the analytic result needs an explicitly justified
interpretation using bounded arithmetic computations and the permitted
induction. Replacing that obligation by the assertion that a classical Lean
proof exists would not prove it.

\section{The logical form of the intended theorem}
\label{sec:logic}

With the fixed predicates \eqref{eq:pos}--\eqref{eq:neg}, define
\begin{align*}
 \Rstr&:\ \forall n>5040\ \exists k\ \Pos(n,k),\\
 \Rdn&:\ \forall n>5040\ \neg\neg\exists k\ \Pos(n,k),\\
 \Rweak&:\ \forall n>5040\ \forall k\ \neg\Neg(n,k).
\end{align*}
The first is the intended positive strict Robin statement. It is
$\Pi^0_2$ in the language with fixed primitive-recursive function
symbols. The last is the absence of a certified strict violation,
a $\Pi^0_1$ formula after pairing its universal variables.

\begin{proposition}\label{prop:logic}
Intuitionistically, $\Rstr\to\Rdn\to\Rweak$.
The weak condition does not imply either of the other two in general.
\end{proposition}
\begin{proof}
The first implication introduces double negation. For the second,
suppose a negative certificate at $n$ is given. Incompatibility of
interval signs refutes every positive certificate at the same $n$,
contradicting $\neg\neg\exists k\,\Pos(n,k)$.

For the failure of a generic reverse rule, use the constant exact
interval $[0,0]$. It has no negative certificate and no positive
certificate, even under double negation. This example concerns a
logical rule; it makes no assertion that an actual Robin gap is zero.
\end{proof}

Replacing the strict inequality by $\Rweak$ thus changes the immediate
constructive content. Negative translation of a classical proof with
$\Rstr$ as conclusion ordinarily yields $\Rdn$, not $\Rstr$.
Removing the remaining double negation would be a Markov-type
witness-extraction step unless a bound or a separate theorem is supplied.
For example, encode a decidable predicate by the rapid rational name
which is zero until its first witness $j$ is found and then equals
$2^{-j}$. Positive apartness of its limit is precisely existence of a
witness. A generic double-negation removal rule would yield Markov's
principle for that predicate.
There is no need to speculate about the irrationality of $e^\gamma$,
and no such speculation would supply the missing proof.

\subsection{A concrete negative RH coding: the eta function}
\label{sec:etacert}

To avoid defining RH through Robin's criterion itself, use certificates
for zeros in rational rectangles strictly to the right of the critical
line. The following specifies a fixed primitive-recursive checker
design. Its analytic adequacy is a separate proof obligation, not a
completed result of the supplied Lean shell.

For $\Re s>0$ use the paired Dirichlet eta series
\[
 \eta(s)=\sum_{j=1}^{\infty}
       \bigl((2j-1)^{-s}-(2j)^{-s}\bigr).
\]
Let $R=[a,b]+i[c,d]$ be rational, with $1/2<a<b<1$, and let the
integer $M\ge1$ bound $|s|$ on $R$, for instance by the sum of the
largest absolute real and imaginary coordinates, rounded upward.
The identity
\[
 (2j-1)^{-s}-(2j)^{-s}
   =s\int_{2j-1}^{2j}t^{-s-1}\,dt
\]
yields the explicit tail
\begin{equation}\label{eq:eta-tail}
 \left|\eta(s)-\sum_{j=1}^{N}
       ((2j-1)^{-s}-(2j)^{-s})\right|
 \le\frac{M}{a}(2N+1)^{-a}
 \le\frac{2M}{\sqrt{2N}}.
\end{equation}
Indeed the omitted disjoint intervals are contained in
$[2N+1,\infty)$, and integration of $t^{-a-1}$ gives the first
inequality. Differentiating the displayed integrand gives
$(1-s\log t)t^{-s-1}$; the integrable majorant proves
\begin{equation}\label{eq:eta-derivative}
 |\eta'(s)|\le \frac1a+\frac{M}{a^2}\le 2+4M.
\end{equation}
These are explicit analytic estimates, with their finite-tail
counterparts available by the same inequalities.

For requested error $\varepsilon=2^{-k}$, choose
$N\ge8M^2\varepsilon^{-2}$. The tail is at most $\varepsilon/2$.
Evaluate each of the $2N$ finite terms with error at most
$\varepsilon/(4N)$. Each term is
$\exp(-s\log j)$, so the rational logarithm algorithm and a bounded
complex exponential Taylor sum suffice. For $|z|\le D$ with integer
$D\ge1$, a particularly crude adequate Taylor cutoff is
\[
 K=2\max(2D^2,p+4).
\]
The factorial lower bound from the final half of the factors and a
geometric ratio bound show that the tail after degree $K$ is at most
$2^{-p-3}$. Here $D$ must bound the approximate argument as well.
For a fully specified, deliberately generous schedule set
\[
 \begin{split}
 N&=8M^2\,2^{2k},\qquad p=k+\ell(N)+4,\\
 D&=M(2N+1)+1,\qquad
 h=p+\ell(M)+2(D+1)+4 .
 \end{split}
\]
For each $j\le2N$, choose the lower endpoint of
$\operatorname{LogI}(j,h)$, use it in $-s\log j$, and evaluate the
complex Taylor polynomial through degree $K$ as above. Since
$\log j\le j$ and the approximation error is at most one, $D$
bounds both arguments. The input perturbation is at most
$3^{D+1}M2^{-h}<2^{-p-4}$, and the Taylor tail is at most
$2^{-p-3}$. Each term therefore has error below
$2^{-p}<\varepsilon/(4N)$. These large but fixed
primitive-recursive bounds specify the rational eta observations
without invoking a universal program evaluator.

\Needspace{6\baselineskip}
The certificate code contains:
\begin{enumerate}[leftmargin=*]
\item The rational rectangle $R$ and an ordered rational subdivision
$z_0,\ldots,z_t=z_0$ of its positively oriented boundary.
\item A precision $k$ and the canonical rational eta approximations
$w_j$ at these vertices, each with coordinate error at most
$\varepsilon$. The checker recomputes them; it does not trust them.
\item A rational ray direction avoiding the polygon vertices, used to
compute an exact winding number.
\end{enumerate}
Put $L=2+4M$ and
\[
 \delta_j=|\Re(z_{j+1}-z_j)|+|\Im(z_{j+1}-z_j)|,\qquad
 E_j=2\varepsilon+L\delta_j.
\]
Require, for each $j$, that the coordinate square centered at $w_j$
of radius $E_j$ avoid zero. This is the rational disjunction
$|\Re w_j|>E_j$ or $|\Im w_j|>E_j$.
The eta image of the boundary segment and the straight polygon chord
from $w_j$ to $w_{j+1}$ both lie in that square. A straight homotopy
therefore avoids zero, and the endpoint homotopies agree between
adjacent segments.

Multiply all values by the complex conjugate of the chosen nonzero
Gaussian-rational direction. This is an orientation-preserving
rational similarity taking that direction to the positive real ray.
Require every transformed vertex to have nonzero imaginary coordinate.
Polygon winding is then computed by signed crossings.
For an edge with endpoint imaginary parts of opposite signs, its
real coordinate at height zero is
\[
 \frac{x_jy_{j+1}-x_{j+1}y_j}{y_{j+1}-y_j}.
\]
An upward crossing on the positive ray contributes $+1$; a downward
one contributes $-1$. Degenerate cases are rejected or avoided by
the supplied direction. Every decision is rational arithmetic.
The checker accepts when all the checks pass and this integer
winding number is positive. Malformed codes return false.
Call this fixed predicate $\Bad(c)$.

The analytic argument principle would establish soundness:
an accepted code encloses a zero of $\eta$. On $b<1$,
\[
 \eta(s)=(1-2^{1-s})\zeta(s),\qquad
 |2^{1-s}|\ge2^{1-b}>1,
\]
so the factor is separated from zero and the enclosed zero is a
zeta zero to the right of the critical line.

Define
\begin{equation}\label{eq:rhc}
 \RHc:\Longleftrightarrow \forall c\ \neg\Bad(c).
\end{equation}
This is a negative $\Pi^0_1$ formula and is intuitionistically stable:
$\neg\neg\RHc\to\RHc$. To see this, fix $c$ with $\Bad(c)$; it
would refute $\RHc$, contrary to the given double negation.

\paragraph{Adequacy that must still be proved.}
The finite evaluator, the homotopy argument, and the argument principle
need certified coded versions. Completeness requires surrounding every
off-line zero with a rational rectangle whose boundary has a positive
lower modulus, then finding an adequate finite subdivision. The functional
equation reflects left-half-strip zeros to the right. If the target RH
definition ranges beyond zeros in $0<\Re s<1$, the zero-free boundary
and exterior, and the trivial zeros, also need the usual interfaces.
None of these internal derivations is supplied by giving the checker
a name. No completed implementation of this checker is claimed here.

In particular the positive assertion that every zero-free compact
rectangle has a finite positive lower-bound certificate must not be
silently substituted for the negative assertion \eqref{eq:rhc}.
That conversion has its own proof requirements.

\subsection{Exactly sufficient certificate transformations}

Suppose four fixed primitive-recursive functions $u,p,v,q$ have been
constructed and $\ipra$ proves
\begin{align}
 n>5040&\ \to\
       \bigl(\Bad(u(n))\vee\Pos(n,p(n))\bigr),\label{eq:F}\\
 \Bad(c)&\ \to\
       \bigl(v(c)>5040\ \wedge\ \Neg(v(c),q(c))\bigr).\label{eq:R}
\end{align}
These statements must concern the actual specified checkers. They are
not axioms adopted in this report.

\begin{theorem}[Conditional assembly]\label{thm:assembly}
If the indicated object proofs of \eqref{eq:F}, \eqref{eq:R}, and
sign incompatibility are supplied, then $\ipra$ proves
$\RHc\leftrightarrow\Rstr$.
\end{theorem}
\begin{proof}
Assume $\RHc$ and fix $n>5040$. It rules out the first disjunct of
\eqref{eq:F}; $p(n)$ is the required positive witness.
Conversely assume $\Rstr$ and a particular $\Bad(c)$.
Equation \eqref{eq:R} gives a negative certificate at $v(c)>5040$,
while $\Rstr$ gives a positive certificate at that same integer.
Their incompatibility gives false.
\end{proof}

The theorem is compiled in the attached Lean file with both
transformations as explicit hypotheses. It is useful proof assembly,
not Robin's theorem with the analytic work concealed.

There is also a less witness-intensive legitimate converse. Because
$\RHc$ is stable, an intuitionistic proof of
$\neg\RHc\to\neg\Rstr$ suffices for $\Rstr\to\RHc$.
One must still prove that contrapositive in the correct theory;
it is not legitimate to turn $\neg\RHc$ into a positive
$\exists c\,\Bad(c)$ by an unstated Markov principle.

\subsection{A positive transfer metatheorem}

\begin{theorem}[Classical PRA is a sufficient intermediate target]
\label{thm:transfer}
Let $A\equiv\forall c\,a(c)$ and
$B\equiv\forall n\,\exists k\,b(n,k)$, where $a,b$ are
quantifier-free formulas over a fixed primitive-recursive language.
If classical first-order $\pra$ proves $A\leftrightarrow B$, then
$\ipra$ proves $A\leftrightarrow B$.
\end{theorem}
\begin{proof}
Classically, $A\to B$ implies the $\Pi^0_2$ sentence
\[
 \forall n\ \exists c,k\quad
       \bigl(\neg a(c)\vee b(n,k)\bigr).
\]
Herbrand witnessing for PRA supplies primitive-recursive terms
$c(n),k(n)$ and a quantifier-free PRA proof of
$\neg a(c(n))\vee b(n,k(n))$; a pair code handles the two witnesses.
The witnessing theorem is stated explicitly in
\cite[p.~13]{AvigadMining}. Quantifier-free arithmetic predicates
are decidable in $\ipra$, so this quantifier-free derivation can be
read intuitionistically. Under $A$ the first disjunct is impossible,
and $k(n)$ supplies the desired witness. This proves $A\to B$.

Apply the standard negative translation to the classical derivation
of $B\to A$. Decidability of quantifier-free formulas identifies
the translation of $A$ with $A$, and the translation of $B$ with
\[
 B^N\equiv\forall n\,\neg\neg\exists k\,b(n,k).
\]
Translated quantifier-free induction instances remain equivalent
to the permitted induction instances. Hence $\ipra$ proves
$B^N\to A$. Since $B\to B^N$ is intuitionistically valid, it also
proves $B\to A$.
\end{proof}

For Robin, fold the decidable condition $n>5040$ into $b(n,k)$ and
use a dummy witness below the cutoff. The theorem then applies to
$\RHc$ and $\Rstr$.

This metatheorem resolves one logical concern: strict witnesses are
not intrinsically incompatible with an initial classical approach.
Its premise is nevertheless a proof in \emph{classical PRA}, not
an arbitrary classical Lean proof. Establishing the necessary
arithmetic strength of the analytic argument remains the central
uncompleted task.

Likewise, conservation of a stronger theory for $\Pi^0_2$ sentences
can be applied to concrete witness statements such as
\eqref{eq:F}--\eqref{eq:R}. It cannot automatically be applied to
an implication with a $\Pi^0_2$ antecedent merely by calling the
whole implication arithmetical.

\section{Where the analytic proof must be completed}
\label{sec:analytic}

This section records the genuine proof obligations behind the classical
criterion. It does not declare them axioms of $\ipra$ or count their
restatement as a proof. Robin's original argument is in
\cite[\S\S3--4]{Robin}; the Nicolas oscillation input is developed in
\cite{Nicolas1983,Nicolas2012}. A modern classical Lean reconstruction
provides a substantially more explicit dependency map \cite{Whidden}.

\subsection{An elementary extremal interpolation lemma}

Let $N,N'$ be integers with $3\le N<N'$. Suppose a real parameter
$\varepsilon>0$ and a constant $C$ give
\[
 \log A(j)\le C+\varepsilon\log j
       \quad(N\le j\le N'),
\]
with equality at both endpoints. Then the Robin logarithmic margin
\[
 M(j)=\log A(j)-\gamma-\log\log\log j
\]
satisfies
\begin{equation}\label{eq:interpolation}
 M(j)\le \max(M(N),M(N'))\qquad(N\le j\le N').
\end{equation}
Indeed, put $x=\log j$ and $f(x)=\varepsilon x-\log\log x$.
On $x>1$,
\[
 f''(x)=\frac{\log x+1}{x^2(\log x)^2}>0.
\]
Thus $f$ is convex and bounded above by the larger of its two endpoint
values on the interval. Adding $C-\gamma$ proves
\eqref{eq:interpolation}. One can replace the maximum statement by a
rational convex-combination upper bound when carrying explicit
certificates.

For colossally abundant endpoints the affine bound comes from maximizing
$A(j)j^{-\varepsilon}$. The displayed lemma itself only needs a
certified bound on the finite interval. Its proof does not require
deciding equality between unrelated transcendental threshold values.
A constructive implementation should allow ties or use overlapping
rational upper bounds, instead of assuming that every threshold
comparison is strictly decidable.

\subsection{The RH-dependent weighted explicit formula}

Put $H=\log n$,
\[
 S(H)=\frac1{\sqrt H\log H},
 \qquad c_0=\gamma+2-\log(4\pi).
\]
The sharp forward proof uses weighted Chebyshev errors
\[
 I_j(x)=\int_x^\infty(\psi(t)-t)t^{-j-1}
                      \frac{j\log t+1}{(\log t)^2}\,dt.
\]
Under RH a multiplicity-counted zero sum yields the estimates
\begin{align*}
 -B_j(x)-\frac{\log(2\pi)x^{-j}}{\log x}
 &\le I_j(x)\le B_j(x),\\
 B_j(x)&=c_0\frac{x^{1/2-j}}{\log x}
       \left(j+\frac1{\log x}
                  +\frac4{(2j-1)(\log x)^2}\right).
\end{align*}
The zero contribution is controlled by
$\sum_\rho|\rho|^{-2}=c_0$ under RH. These statements involve
infinite zero sums and integrals, not only evaluations of elementary
constants. The required arithmetic version must state finite contour
identities, prove their error estimates, control the zero multiplicities,
and provide bounds for discarded tails.

For example, a formal statement that a sequence of approximants converges
does not provide the modulus needed to replace this integral by a fixed
primitive-recursive observation. Likewise a type enumerating zeros is
not a proved bound on the tail of their reciprocal squares. The
quantitative packets already present in ProveIt are relevant candidate
arguments, but their stated proof status must be respected.

\subsection{A concrete sharp margin available after that work}

The classical reconstruction \cite{Whidden} organizes the forward estimate
as
\begin{equation}\label{eq:budget}
 M(n)\le E(H)-C(H)\qquad(H\ge20000,\ \RH),
\end{equation}
where
\begin{align*}
 C(H)&=\sum_{p\le H/2}
       \min\left(\frac{\log p}{H\log H},\frac1{p^2}\right),\\
 E(H)&=(2+c_0)\frac{H^{-1/2}}{\log H}
       +(c_0-2)\frac{H^{-1/2}}{(\log H)^2}\\
 &\quad +(8+4c_0)\frac{H^{-1/2}}{(\log H)^3}
       +2\frac{H^{-2/3}}{\log H}
       +\log(2\pi)\frac{H^{-1}}{\log H}.
\end{align*}
The first term is an explicit sum of losses forced by prime powers;
the second is an upper error estimate obtained from the weighted
explicit formula. The published reconstruction gives, for $H\ge74500$,
\begin{equation}\label{eq:rational-budgets}
 \frac{E(H)}{S(H)}
 \le\frac{3010451}{1332800}
 =\frac{113}{50}-\frac{1677}{1332800},
 \qquad
 \frac{C(H)}{S(H)}\ge\frac{113}{50}.
\end{equation}
These are recorded as concrete porting targets. This report has not
reproved their analytic premises in $\ipra$, nor independently rebuilt
the external Lean development.

Their implication for constructive strictness can, however, be derived
explicitly.

\begin{proposition}[Conditional positive separation]\label{prop:separation}
Assume the actual estimates \eqref{eq:budget} and
\eqref{eq:rational-budgets} for an integer $n$ with $\log n\ge74500$.
Let $a=1677/1332800$. Then
\[
 B(n)-A(n)\ge\frac{a}{2n^2}
             =\frac{1677}{2665600\,n^2}.
\]
Consequently the canonical interval test $\Pos(n,k)$ succeeds at
$k=2\ell(n)+12$.
\end{proposition}
\begin{proof}
The estimates imply $M(n)\le-aS(H)$.
For $n$ in this range, $\sqrt H\le n$ and $\log H\le n$, so
$S(H)\ge n^{-2}$ and $M(n)\le-a/n^2$.
Since $M(n)=\log(A(n)/B(n))$ and $B(n)\ge1$,
\[
 B(n)-A(n)\ge B(n)(1-e^{-a/n^2}).
\]
For $0\le u\le1$, $e^u\ge1+u$, and therefore
$1-e^{-u}\ge u/(1+u)\ge u/2$. Use $u=a/n^2$.

Finally $2^{-2\ell(n)-12}\le1/(4096n^2)<a/(2n^2)$.
An interval for $B(n)$ with width at most this quantity has lower
endpoint strictly above $A(n)$. The test \eqref{eq:pos} therefore
accepts at the stated primitive-recursive precision.
\end{proof}

This proposition shows why the strict inequality should not be treated
as a permanent obstruction. The forward argument already contains
room for a useful explicit positive witness. The unresolved question
is the justification of its analytic input in the chosen arithmetic
theory.

\subsection{The finite startup range is another real obligation}

The same external project reports three stages: exact factorizations
for $5041\le n<7560$, a common tangent up to $720720$, and $36$
rational logarithmic-height rows reaching the analytic range.
Its README reports ordinary kernel checks of the finite data and an
axiom audit using only \code{propext}, \code{Classical.choice}, and
\code{Quot.sound}. Those are reports of the external project, not
checks rerun in this task.

The finite verifier and its soundness proof are plausible objects to
port because the accepted rows are rational finite data. But a
single finite bound can be enormous, and a classical real-valued
soundness theorem still needs translation. The present archive's
verification through $100000$ does not cover this logarithmic-height
range and is not offered as a replacement for it.

A completed constructive argument should make the handoff using a
decidable integer cutoff, or overlapping rational height ranges.
It should not invoke the general real trichotomy
$\log n<74500$ or $\log n\ge74500$ without explaining its arithmetic
realization. The existing rational rows provide a useful approach to
that engineering problem, but no such port is claimed here.

\subsection{What the converse still needs}

The classical Nicolas--Landau mechanism supplies negative excursions
of $\log F(x)$ when an off-critical zero exists. In a convenient form,
there are $c>0$ and $0<b<1/2$ such that arbitrarily far out one has
\[
 -\log F(Q)\ge cQ^{-b}.
\]
Because $\vartheta(x)$ and the prime product are constant between
successive integers, replacing a real frontier by its integer part
only weakens the constant by a fixed factor. Thus an integer-frontier
form can be used.

The preceding LCM section proves the remaining finite implication
in full detail. Once an excursion beyond its explicit rational
cutoff is supplied, its output is an actual strict counterexample.
There is no additional number-theoretic asymptotic argument hidden
in that transfer.

The missing analytic proof is substantial. Its familiar ingredients
include analytic continuation of a weighted prime-function transform,
Landau's positive-transform principle, and the detection of a
singularity associated with an off-line zero. A claim of arbitrary
large excursions in classical analysis is not yet a
primitive-recursive witness bound in $\ipra$.

There are two defensible completion routes:
\begin{enumerate}[leftmargin=*]
\item Starting with a specified accepted zero certificate, prove a
quantitative transform theorem producing a frontier, a rational
excursion gap, and a primitive-recursive bound for their search.
The LCM transfer then gives \eqref{eq:R}.
\item Under $\Rstr$ and an accepted zero certificate, derive an
analytic contradiction inside a proved conservative arithmetic
framework. For the negative RH conclusion, a direct contradiction
can suffice; no positive counterexample selector need be asserted
unless it is actually obtained.
\end{enumerate}

The second route may be shorter. Neither is completed by citing
ordinary Lean's classical choice or by declaring the needed
oscillation theorem as a structure field.

\subsection{How conservation results could help, and their limit}

Sudac's work \cite{Sudac} develops real and complex analysis in
extensions conservative over PRA and obtains a PRA proof of the
prime number theorem. It is evidence that substantial analytic
number theory can admit such treatment. It is not a published
PRA proof of Robin's equivalence, and no implication from PNT
alone gives the sharp RH-sensitive estimates here.

A serious alternative to direct object-proof construction is to
check every needed analytic principle in an appropriate conservative
theory and formalize the finite witness statements there.
The hypothesis restrictions and the exact conservation class must
be verified. The transfer theorem in Section~\ref{sec:logic}
explains why a genuine classical PRA endpoint would suffice
afterward. This is a research route with a concrete endpoint,
not a substitute for the missing derivation.

\section{Lean formalization and reuse of ProveIt}
\label{sec:lean}

\subsection{The locally checked Lean component}

The file \code{lean/RobinCertificateLogic.lean} imports only
\code{Std}. Its definitions have the following structure:
\begin{lstlisting}
structure Checkers where
  positive : Nat -> Nat -> Bool
  negative : Nat -> Nat -> Bool
  badZero : Nat -> Bool

def RobinStrict (c : Checkers) : Prop :=
  forall n, 5040 < n ->
    Exists (fun k => c.positive n k = true)

def RHNegative (c : Checkers) : Prop :=
  forall z, Not (c.badZero z = true)
\end{lstlisting}
The shell defines the weak and double-negated forms as well.
Its assembly theorem,
\begin{center}
\code{criterion\_of\_certificate\_transforms},
\end{center}
takes four
selector functions, interval sign incompatibility, and the two
transformations \eqref{eq:F}--\eqref{eq:R} as explicit parameters.
Its constructive proof is exactly Theorem~\ref{thm:assembly}.

The file also implements bounded search:
\begin{align*}
 \operatorname{any}(p,0)&=p(0),\\
 \operatorname{any}(p,b+1)&=\operatorname{any}(p,b)\vee p(b+1),\\
 \operatorname{pick}(p,0)&=0,\\
 \operatorname{pick}(p,b+1)&=
 \begin{cases}
 b+1,&p(b+1)=\mathsf{true},\\
 \operatorname{pick}(p,b),&p(b+1)=\mathsf{false}.
 \end{cases}
\end{align*}
It proves the two invariants
\[
 \operatorname{pick}(p,b)\le b,\qquad
 \operatorname{any}(p,b)=\mathsf{true}
 \ \to\ p(\operatorname{pick}(p,b))=\mathsf{true}.
\]
After instantiating $p$ with one fixed PR predicate, these are
quantifier-free induction invariants. In the successor step, the
decidable Boolean value $p(b+1)$ either gives the new witness or
reduces the assertion to the previous invariant.

The locally used compiler was Lean 4.32.0. The following eight
declarations were audited with \code{\#print axioms}:
\begin{center}
\begin{tabular}{@{}l@{}}
\toprule
\code{strict\_implies\_doubleNeg}\\
\code{doubleNeg\_implies\_weak}\\
\code{rhNegative\_stable}\\
\code{converse\_of\_contrapositive}\\
\code{criterion\_of\_certificate\_transforms}\\
\code{zeroGap\_is\_not\_strict}\\
\code{pickUpTo\_le}\\
\code{pickUpTo\_correct}\\
\bottomrule
\end{tabular}
\end{center}
All eight reported no axioms. The file contains no admitted proof,
custom axiom declaration, classical choice, or native evaluation
proof shortcut. This is a check of the logical shell, not an
analytic Robin proof. Even an axiom-free ambient Lean theorem
need not, merely by that fact, be an iPRA theorem; the specific
logical and bounded-induction arguments here are what make the
translation straightforward.

\subsection{What the repository actually provides}

The ProveIt inspection was pinned at commit
\[
 \texttt{\small eb8eb6653c85aa44f1e18dd4efa243f41b5a9e98}.
\]
The root iPRA README separates the checked proof infrastructure
from delivered research packets which have not been verified as
packets \cite{ProveIt}. The relevant full source files and their
immutable URLs are included under \code{sources/proveit/}.
These files were read; the complete repository was not rebuilt
in this task.

\begin{longtable}{@{}p{.33\textwidth}p{.59\textwidth}@{}}
\toprule
Source & Reusable content and boundary\\
\midrule
\endhead
\code{PR.lean} &
The actual PR constructors, term syntax, nonlogical axioms and
quantifier-free induction of the object theory.\\
\code{NativeBiInterpretation.lean} &
A proved equivalence between an existing PRA derivation and its
native presentation. It transports a proof; it does not supply the
analytic theorem.\\
\code{ExponentialTail.lean} &
Finite, denominator-cleared exponential-tail derivations, with
an explicit index $2B^2+2B+p+4$. Analytic identification of the
infinite function is outside that module's stated boundary.\\
\code{FiniteLogarithm.lean} &
Finite alternating logarithm polynomials, exact coefficient
relations, and a derivative residual. The general analytic
identification remains separate.\\
\code{Harmonic.lean} &
A finite arithmetic derivation of $H_{2^{2t}}\ge1+t$.
It does not itself provide the Euler-constant enclosure proved here.\\
Zeta functional-equation and general zeta packets &
Quantitative human arguments and Lean outlines. Their READMEs
explicitly do not claim completed analytic object derivations.\\
Zeta-zero \code{ObjectTargets.lean} &
A \code{Completed} structure requesting contour, nonvanishing,
symmetry, and tail proofs. No inhabitant supplying those proofs
is constructed in the inspected file.\\
RH \code{PAStatement} research note &
An arithmetic statement project using a Mertens/Littlewood
criterion. The semantics and equivalence with analytic RH
are explicitly future work.\\
\bottomrule
\end{longtable}

The exact existing transfer endpoint is
\begin{lstlisting}
(Gamma |- [PRAAx] phi) <->
  NPrf (Gamma.map NFm.ofPRA) (NFm.ofPRA phi)
\end{lstlisting}
where the display uses ASCII for readability. The actual theorem
is \code{pra\_iff\_native}. It requires an actual proof on one
side. Supplying an uninhabited completion record cannot invoke
it to establish a new theorem.

\subsection{An explicit implementation plan}

The following is a dependency plan for completion. The module names
are proposed names, not claims of files already built.

\begin{enumerate}[leftmargin=*]
\item \textbf{Arithmetic encoding.}
Instantiate the fixed PR functions for rational arithmetic,
factorization, $\sigma$, $L_P$, bounded products, interval comparison,
and bit length in the object syntax. Reuse factor-trace and
bounded-syntax modules where their exact statements fit.
Prove each loop invariant by a quantifier-free induction instance.

\item \textbf{Canonical elementary names.}
Formalize \eqref{eq:logtail}, \eqref{eq:gamma-bound},
\eqref{eq:exptail}, and the composition bounds.
For log addition and exponential multiplication, use the finite
binomial/convolution identities and the explicit tail bounds:
compare both sides at arbitrary supplied precision, bound the
finite truncation differences, and clear denominators.
Do not infer an object identity just from an ambient real identity.
This creates fixed terms representing $\Pos$ and $\Neg$ and proves
their incompatibility.

\item \textbf{Finite Robin certificates.}
Compile the dyadic verification recurrence and its endpoint
soundness. Replay the included finite range as a closed arithmetic
calculation, then port the external tangent-row method and its
large finite startup range. The present CSV alone is insufficient
for the analytic handoff.

\item \textbf{LCM transfer.}
Represent $\vartheta$ and $\psi$ as finite sums of log names.
Prove the finite product identity, the binomial divisibility
lemma, and the $8/m$ and $4/m$ estimates with supplied precision.
An input rational excursion gap determines the accuracy needed
at every use of the identity. This yields a proved transformer
from a certified excursion to $\Neg(N,k)$.

\item \textbf{Eta and zero certificates.}
Implement the exact evaluator and the rational polygon checker
specified in Section~\ref{sec:etacert}. Internalize their soundness,
argument-principle meaning, completeness under appropriate zero
data, and reflection. Identify the precise ordinary RH statement
to which this coding is adequate.

\item \textbf{Weighted explicit formula.}
Prove finite-height contour identities with rational error budgets;
separate poles, trivial zeros and counted nontrivial zeros; prove
the reciprocal-zero tails and the weighted Chebyshev estimates.
This is a central analytic task, not a routine wrapper lemma.
Audit whether the negative RH formulation supplies each needed
positive lower bound or whether the proof must be reorganized.

\item \textbf{Forward estimate and strict margin.}
Internalize the exponent budget \eqref{eq:budget} and the rational
coefficient inequalities \eqref{eq:rational-budgets}. Apply
Proposition~\ref{prop:separation} to get an explicit precision
for every large integer. Combine with the finite startup proof
using an arithmetic cutoff.

\item \textbf{Converse analytic input.}
Internalize the Nicolas--Landau argument through a quantitative
certificate transformer or a direct stable contradiction,
as distinguished in Section~\ref{sec:analytic}. The already proved
finite LCM reduction handles the final arithmetic step.

\item \textbf{Object theorem and presentation transfer.}
Construct an actual value of
\[
 \operatorname{Prf}\ \operatorname{PRAAx}\ []
              \ \operatorname{RobinEquivalenceSentence}.
\]
Only then apply \code{pra\_iff\_native}. If a conservative classical
route is used instead, first obtain the actual classical PRA
derivation and then implement the proof transformation of
Theorem~\ref{thm:transfer}.
\end{enumerate}

For the intended object sentence, $\RHc$ is a universal closure
of the negation of the Boolean zero-check term, while $\Rstr$ is
the universal $n$-closure of a decidable cutoff implication and
the existential $k$-closure of the positive-check term.
The equivalence is their conjunction of two implications in the
existing intuitionistic formula syntax. Thus the endpoint is
specified without quantifying over a Lean $\RR$ or $\CC$.

No object implementation of steps 5--8 is supplied here. They are
the substantive outstanding work, and the roadmap does not assert
that their feasibility has already been proved.

\section{Reproduction and evidence}

From the extracted archive root, the exact finite replay is
\begin{lstlisting}
python3 code/robin_certificates.py \
  --verify certificates/finite_robin_certificates.csv
\end{lstlisting}
\Needspace{6\baselineskip}
To regenerate the certificate table:
\begin{lstlisting}
python3 code/robin_certificates.py \
  --limit 100000 --output certificates
\end{lstlisting}
The independent arithmetic audit is
\begin{lstlisting}
python3 verification/independent_arithmetic_audit.py
\end{lstlisting}
It uses exact \code{Fraction} references rather than the production
dyadic recurrences. It checked $141$ logarithm cases, an independent
$m=8192$ Euler-constant enclosure, seven exponential cases, $49$
signed products and $19$ composed Robin bounds. It also recomputed
all $99999$ divisor sums by prime-factor geometric products and
rechecked their signs. The stored JSON records all checks passing.

With Lean 4.32.0 installed, compile from the \code{lean/} directory:
\begin{lstlisting}
lean RobinCertificateLogic.lean
\end{lstlisting}
The companion \code{lean-toolchain} pins that compiler. The archive
does not require Mathlib for this checked component.
The full external classical project requires a different pinned
toolchain and dependencies; it was inspected as source, not rebuilt.

Rebuild the report from \code{paper/}:
\begin{lstlisting}
pdflatex -halt-on-error report.tex
pdflatex -halt-on-error report.tex
\end{lstlisting}
No external bibliography processor is needed. The delivered PDF
was compiled and visually inspected. A file inventory with
SHA-256 digests is included for the substantive archive contents.

\paragraph{Final status.}
The archive provides reusable constructive mathematics, exact
finite evidence and a checked logical core. It does not contain
a completed $\ipra$ proof of Robin's equivalence. In particular
the absent analytic derivations have not been replaced by axioms,
unverified completeness fields, or an assertion that computability
alone establishes primitive-recursive provability.

\begin{thebibliography}{99}
\raggedright

\bibitem{MathWorld}
E. W. Weisstein, \emph{Robin's Theorem}, MathWorld.
\url{https://mathworld.wolfram.com/RobinsTheorem.html}.
Statement checked 19 November 2026.

\bibitem{Robin}
G. Robin, \emph{Grandes valeurs de la fonction somme des diviseurs et
hypoth\`ese de Riemann}, Journal de Math\'ematiques Pures et
Appliqu\'ees \textbf{63} (1984), 187--213.
Theorem 1, p.~188; \S3, Proposition 1, p.~192;
\S3, Proposition 3, p.~204; \S4, pp.~204--205.
Online scan/OCR consulted:
\url{https://fr.scribd.com/document/812768752/Robin-1984-noBackground}.
OCR formulas require checking against the mathematics.

\bibitem{Nicolas1983}
J.-L. Nicolas, \emph{Petites valeurs de la fonction d'Euler},
Journal of Number Theory \textbf{17} (1983), 375--388.
\url{https://doi.org/10.1016/0022-314X(83)90055-0}.

\bibitem{Nicolas2012}
J.-L. Nicolas, \emph{Small values of the Euler function and the
Riemann hypothesis}, Acta Arithmetica \textbf{155} (2012), 311--321.
\url{https://arxiv.org/abs/1202.0729}.

\bibitem{Whidden}
J. Whidden, \emph{Robin's 1984 criterion for the Riemann hypothesis:
a formally verified proof}, September 2026, companion to the
\code{Robin1984} Lean project.
\url{https://github.com/kimihiro64/Robin1984}.
Observed commit:
\texttt{acab1a31f31e0499518a4416b63281e0b4838f9c}.
Paper path: \path{paper/robin1984-formalization.tex};
observed paper blob:
\texttt{0971f7fdf590715a5d91c894015a996a3461bc13}.
README and source inspected; project build and reported independent
kernel audit not replayed here.

\bibitem{AvigadMining}
J. Avigad, \emph{Proof mining}, ASL 2004 lectures, especially
p.~13 on Herbrand witnessing for PRA and pp.~16--17 on
restricted versus unrestricted higher-type recursion.
\url{https://www.andrew.cmu.edu/user/avigad/Talks/asl04.pdf}.

\bibitem{Sudac}
O. Sudac, \emph{The prime number theorem is PRA-provable},
Theoretical Computer Science \textbf{257} (2001), 185--239.
\url{https://doi.org/10.1016/S0304-3975(00)00116-X}.
The published abstract establishes the stated scope of the work;
no Robin-specific theorem is attributed to it.

\bibitem{ProveIt}
V. Reshetnikov et al., \emph{ProveIt}, source snapshot
\texttt{eb8eb6653c85aa44f1e18dd4efa243f41b5a9e98}.
\url{https://github.com/VladimirReshetnikov/ProveIt}.
Relevant tree: \path{Logic/PrimitiveRecursiveArithmetic}.
Exact paths, Git blob hashes and immutable URLs for the inspected
files are supplied in \path{sources/proveit/manifest.json};
the detailed audit is \path{sources/proveit_audit.md}.
The RH statement-only research note is under
\path{NumberTheory/RiemannHypothesis/PAStatement/Research/}.

\end{thebibliography}

\end{document}
